% ============================================================
% I. INTRODUCTION
% ============================================================
\section{Introduction}
\label{sec:intro}

% ---- 1. Problem context ----
Robot simulation, motion planning, and control depend on efficient collision geometry that is lightweight enough for real-time queries yet faithful enough to avoid missed collisions.
In the ROS ecosystem, URDF~\cite{urdf_doc} supports both analytic primitives (spheres, cylinders, boxes) and triangular meshes.
In practice, robot descriptions carry high-detail visual meshes alongside simplified or manually authored collision approximations, yielding representations that are either computationally expensive or inconsistently maintained.

% ---- 2. Existing practice -%
Current collision geometry authoring is largely manual: designers hand-place primitives, simplify meshes with modeling tools, or rely on convex decomposition.
These approaches are time-consuming, hard to reproduce, and produce collision models at arbitrary fidelity levels.
Simulators and planners---MuJoCo~\cite{todorov2012mujoco}, Drake~\cite{drake}, PyBullet~\cite{coumans2018pybullet}---support analytic collision geometries with closed-form distance computation, but generating such primitives from existing URDF descriptions within an open-source, URDF-native toolchain has not been addressed.

% ---- 3. Analytic primitives -%
Analytic primitives (spheres, capsules, boxes) offer significant advantages over meshes: stable signed-distance evaluation, analytical contact gradients, and efficient broad-phase culling~\cite{pan2012fcl}.
Capsules (swept spheres) match the elongated shape of robot links, and their distance function---point-to-segment distance minus radius---evaluates in constant time.
Despite these advantages, converting URDF mesh geometry into analytic capsule primitives with measured vertex coverage remains an open tooling gap.

% ---- 4. Our approach -%
In this paper we introduce URDFApproxGeom, an integrated open-source toolchain that reads a URDF with mesh collision (or visual) geometry and outputs drop-in replacement URDFs with convex, sphere-tree, or capsule collision approximations.
The toolchain is available as a Python package with a CLI, a public Python API, Docker support, and bundled robot resources.
Our primary methodological contribution is an automatic capsule-fitting pipeline that proceeds in five stages:
(i)~computing the principal axis of each link mesh via PCA;
(ii)~slicing the mesh with planes perpendicular to this axis and extracting 2D cross-section contours;
(iii)~fitting covering circles to each cross section using minimum enclosing circles and, where needed, adaptive Lloyd-style refinement to handle concave or multi-lobed sections;
(iv)~chaining matched circles across adjacent sections into capsules, followed by collinear merging, vertex-coverage-aware radius growth, endpoint tightening, inflation-based splitting, and nested-capsule pruning;
(v)~emitting each capsule as a URDF-compatible cylinder-plus-two-sphere decomposition alongside a JSON sidecar with the canonical analytic parameters.

% ---- 5. Broader system -%
Beyond capsule fitting, the toolchain provides two additional approximation modes for completeness: convex hull generation via CGAL/libigl and sphere-tree approximation via a medial-axis tree algorithm.
All modes emit standard URDF and are accompanied by a CLI validation utility that reports coverage, volume tightness, and radius-inflation metrics.
A built-in comparison workflow generates side-by-side visualizations across all modes and presets, enabling qualitative and quantitative inspection without external tools.

% ---- 6. Contributions -%
In summary, this paper contributes:
\begin{itemize}
    \item \textbf{An integrated URDF-to-primitive approximation pipeline} that preserves visual, inertial, and joint structure while replacing collision geometry with convex, sphere, or capsule primitives.
    \item \textbf{A capsule-output pipeline} that converts robot link meshes into low-count analytic capsule primitives through PCA sectioning, adaptive circle fitting, radius growth, and tightness-guided refinement, with coverage treated as an explicit validation metric rather than an assumed guarantee.
    \item \textbf{A reproducible CLI/Python workflow} with named presets, validation metrics, pass/fail gates, and visualization bundles for qualitative comparison.
    \item \textbf{An artifact-backed empirical characterization} on the Franka FR3 robot, reporting primitive count, runtime, vertex coverage error, capsule union volume, and preset sensitivity across all three modes.
\end{itemize}
